\ifx\danslelivre\undefined
\documentclass[10pt]{../fichiers-configuration-latex/smfart}
\usepackage{ifthen}
\newboolean{english}
\newboolean{francais}
\setboolean{english}{false}
\setboolean{francais}{true}
\input{../fichiers-configuration-latex/paquets}
\selectlanguage{francais}
\input{../fichiers-configuration-latex/lemme-proposition-theoreme}
\input{../fichiers-configuration-latex/numerotation}
\input{../fichiers-configuration-latex/macros}

\externaldocument{../01-quasi-excellents/excellent}
\externaldocument{../02-topologies/topologies}
\externaldocument{../03-approximation/approximation}
\externaldocument{../04-Cohen-Gabber/Cohen-Gabber}
\externaldocument{../05-algebrisation-partielle/algebrisation}
\externaldocument{../06-log-lissite/log}
\externaldocument{../07-uniformisation-locale/uniformisation-faible}
%\externaldocument{../08-modification-theorem/modification}
\externaldocument{../09-uniformisation-premiere-a-l/uniformisation-locale-premiere-a-l}
\externaldocument{../11-produits-orientes/produits-orientes}
\externaldocument{../12-descente-cohomologique-orientee/descente-orientee}
\externaldocument{../13-theoreme-de-finitude/finitude}
\externaldocument{../14-fonctions-dimension/fonctions-dimension}
\externaldocument{../15-Lefschetz-affine/Lefschetz-affine}
\externaldocument{../16-classes-de-Chern/gysin}
\externaldocument{../17-complexes-dualisants/dualite}
\externaldocument{../20-rigidite/rigidite}
\externaldocument{../21-finitude-non-abelienne/finitude-non-abelienne}

\author{Luc Illusie}

\title{Introduction}

\begin{document}

\maketitle

\else
\setboolean{english}{true}
\setboolean{francais}{false}
\selectlanguage{francais}
\chapter{Introduction}
\addtocontents{toc}{par Luc Illusie}
\begin{center}
Luc Illusie
\end{center}
\fi

\renewcommand\thesubsection{\textbf{\arabic{subsection}}}

Le présent volume rassemble les exposés d'un groupe de travail qui s'est tenu à l'École
polytechnique du printemps 2006 au printemps 2008 sur les travaux de Gabber présentés dans
\cite{Conference-Illusie@Gabber} et \cite{Conference-Deligne@Gabber}. Ceux-ci portent sur la
cohomologie étale et l'uniformisation des schémas quasi-excellents. 

En ce qui concerne la cohomologie étale, un résultat central est le théorème de finitude suivant :

\begin{theoreme}\label{0-finitude1} Soient $Y$ un schéma noethérien quasi-excellent (\textit{cf.}
\refext{I}{1def-excellent}), $f : X \ra Y$ un morphisme de type fini, $n$ un entier inversible sur
$Y$, et $F$ un faisceau constructible de $\ZZ/n\ZZ$-modules sur X. Alors, pour tout $q$, $R^qf_*F$
est constructible, et il existe un entier $N$ tel que $R^qf_*F = 0$ pour $q \ge N$.  \end{theoreme} 

Rappelons que, sans hypothèse sur $Y$ ni sur $n$, mais lorsqu'on suppose $f$ propre et de
présentation finie, les faisceaux $R^qf_*F$ sont constructibles, et nuls pour $q > 2d$ si $d$ majore
la dimension des fibres de $f$ \sga{4}{xiv}{1.1}. Si $f$ n'est pas supposé propre, l'hypothèse que
$n$ soit inversible sur $Y$ est essentielle : si $k$ est un corps algébriquement clos de
caractéristique $p >0$ et $X$ la droite affine sur $Y = \SP(k)$, $H^1(X,\ZZ/p\ZZ)$ est un
$\FF_p$-espace vectoriel de dimension infinie (\textit{cf.} \sga{4}{xiv}{1.3}). Pour $Y$ excellent
de caractéristique nulle, \ref{0-finitude1} est démontré dans l'exposé d'Artin \sgacourt{4}{xix}{}. Si
$S$ est un schéma noethérien régulier de dimension $\le 1$ (non nécessairement quasi-excellent) et
$f$ un morphisme de $S$-schémas de type fini,  la conclusion de \ref{0-finitude1} est encore vraie,
d'après \sgaquatreetdemi{Th. Finitude}{1.1} (ce résultat peut d'ailleurs se déduire de
\ref{0-finitude1}, \textit{cf.} XIII).  

La démonstration d'Artin dans \sgacourt{4}{xix}{} utilise la résolution des singularités de Hironaka pour
se ramener au cas où $f$ est l'inclusion du complément d'un diviseur régulier dans un schéma
régulier et $F$ un faisceau constant, auquel cas la conclusion découle du théorème de pureté
cohomologique absolu établi également dans (\textit{loc. cit.}). 

La démonstration de Gabber de \ref{0-finitude1} suit la même méthode, mais :

(a) on doit faire appel au théorème de pureté cohomologique absolu établi dans le cas général par
Gabber \cite{AZUMINO@Fujiwara}, 

(b) on ne dispose plus de la résolution des singularités sous la forme de Hironaka ; celle-ci est
remplacée par un théorème d'uniformisation locale, dû à Gabber (\cite{Conference-Deligne@Gabber}),
qui s'énonce ainsi (\refext{II}{2-théorème uniformisation faible}, \refext{III}{3-unif},
\refext{IX}{9-unif}) : \begin{theoreme}\label{0-uniformisation l'} Soient $X$ un schéma noethérien
quasi-excellent, $Z$ un fermé rare de $X$ et $\ell$ un nombre premier inversible sur $X$. Il existe
une famille finie de morphismes $(p_i : X_i \ra X)_{i \in I}$, couvrante pour la topologie des
$\ell'$-altérations et telle que, pour tout $i \in I$ :

(i) $X_i$ soit régulier et connexe,

(ii) $p_i^{-1}(Z)$ soit le support d'un diviseur à croisements normaux stricts.  \end{theoreme} La
topologie des $\ell'$-altérations est une topologie du type de celles considérées par Voevodsky
(\textit{cf.} \cite{Cohomological@Goodwillie-Lichtenbaum}, pour laquelle les $\ell'$-altérations
(i.~e. les morphismes propres surjectifs génériquement finis de degré résiduel générique premier à
$\ell$) et les recouvrements étales complètement décomposés (i.~e.  \textit{de Nisnevich}) sont des
familles couvrantes, voir (\refext{II}{2-alt}) pour une définition précise.

La première partie de ce volume est consacrée, après des rappels, dans I, sur les notions de schéma
quasi-excellent ou excellent, à la démonstration de \ref{0-uniformisation l'} et de certains
compléments et variantes. 

Il y a trois grandes étapes.

(A) \textit{Fibration en courbes}. La question étant locale pour la topologie de Nisnevich, et en
particulier pour la topologie de Zariski, on peut supposer $X$ de dimension finie. On raisonne par
récurrence sur la dimension $d$ de $X$. On peut supposer $X$ local hensélien. L'hypothèse de
quasi-excellence sur $X$ permet d'appliquer le théorème d'algébrisation de Popescu au complété
$\widehat X$ de $X$, et de se ramener, via des techniques d'approximation dues à Gabber, expliquées
dans III, au cas où $X$ est local complet de dimension $d$, et même intègre, normal, avec $d \ge 2$.
On dispose alors des classiques théorèmes de structure de Cohen. Tels quels, ils sont toutefois
insuffisants. Mais un délicat raffinement, dû à Gabber, démontré dans IV, permet de se ramener,
quitte à remplacer $X$ par une extension finie de degré générique premier à $\ell$, au cas où $X$
est le complété en un point fermé d'un schéma $X'$ de type fini et de dimension relative 1 sur un
schéma local noethérien régulier complet de dimension $d-1$, et le fermé $Z$ le complété d'un fermé
rare $Z'$ de $X'$. Ce théorème de « fibration », ou « d'algébrisation partielle », est
établi dans V. Après quelques nouvelles réductions faciles, on est ramené au cas où $X$ est un
schéma normal intègre, propre sur un schéma affine normal intègre excellent $Y$ de dimension $d-1$,
à fibre générique géométriquement intègre, lisse, et de dimension 1, et le fermé $Z$ un diviseur de
fibre générique étale. 

(B) \textit{de Jong et log régularité}. On peut alors appliquer le théorème de la courbe nodale de
de Jong, sous sa forme équivariante, \cite[2.4]{Families@de_Jong} : il existe un groupe fini $G$ et
une altération projective $G$-équivariante du morphisme $f : X \ra Y$ en un morphisme $f' : X' \ra
Y'$, qui est une courbe projective nodale $G$-équivariante, et l'image inverse $Z'$ de $Z$ un
diviseur de composantes dominantes étales et transverses au lieu lisse de $f'$ (\refext{IX}{9-dJ}).
Appliquant l'hypothèse de récurrence au quotient de $Y'$ par un $\ell$-Sylow de $G$, et utilisant le
lemme d'Abhyankar, on se ramène finalement au cas où $G$ est un $\ell$-groupe, $X = X'/G$, $Y =
Y'/G$, $Y'$ contient un fermé ($G$-équivariant) $T'$ tel que le couple $(Y',T')$ soit \textit{log
régulier}, ainsi que le couple $(X',f'^{-1}(T') \cup D)$, où $D$ est un diviseur ($G$-équivariant)
étale sur $Y'$, avec $Z' \subset f'^{-1}(T') \cup D$ (\textit{cf.} \refext{VI}{6-nodal implique log
lisse}, \refext{VI}{6-Kummer}). 

(C) \textit{Modifications équivariantes}. Le quotient d'un log schéma log régulier par l'action d'un
groupe fini n'est pas en général log régulier. En particulier, $X = X'/G$, muni du fermé
$(f'^{-1}(T') \cup D)/G$, n'est pas nécessairement log régulier. Si ce couple était log régulier, la
désingularisation de Kato-Niziol des schémas log réguliers (\cite{Toric@Kato}, \cite{Toric@Niziol}),
généralisant la classique désingularisation des variétés toriques \cite{Toroidal@KKMS}, terminerait
la démonstration de \ref{0-uniformisation l'}. Gabber a dégagé des conditions suffisantes assurant
que le quotient d'un log schéma log régulier par un groupe fini est encore log régulier. Il s'agit
de la notion d'\textit{action très modérée}, étudiée dans VI.  Gabber montre qu'il existe une
modification projective $G$-équivariante $p : X'' \ra X'$ et un fermé $G$-équivariant $D''$ de $X''$
contenant l'image inverse de $f'^{-1}(T') \cup D$ tels que le couple $(X'',D'')$ soit log régulier,
et l'action de $G$ sur $(X'',D'')$ très modérée. On conclut alors en appliquant la désingularisation
de Kato-Niziol au quotient de $(X'',D'')$ par $G$. La démonstration de ce théorème de modification
est donnée en VIII. Elle s'appuie sur la théorie des désingularisations canoniques en
caractéristique nulle (Hironaka, Bierstone-Milman, Villamayor, Temkin). 

Gabber établit dans X une \textit{variante relative} de ce théorème de modification, où le log
schéma $G$-équivariant considéré est non seulement log régulier, mais log lisse sur une base log
régulière $S$, avec action triviale de $G$ : on peut construire une modification équivariante
respectant la log lissité sur $S$. Il en déduit notamment les raffinements suivants de théorèmes de
de Jong :
\begin{theoreme}\label{0-deJong l'} (1) Soit $X$ un schéma séparé et de type fini sur un
corps $k$, $Z \subset X$ un fermé rare, $\ell$ un nombre premier $\ne \mathop{\rm car}(k)$. Il
existe alors une extension finie $k'$ de $k$ de degré premier à $\ell$ et une $\ell'$-altération $h
: X' \rightarrow X$ au-dessus de $\mathop{\rm Spec}k' \rightarrow \mathop{\rm Spec}k$, avec $X'$
lisse et quasi-projectif sur $k'$, et $h^{-1}(Z)$ le support d'un diviseur à croisements normaux
stricts.

(2) Soient $S$ un schéma noethérien séparé, intègre, excellent, régulier,  de dimension 1, de point
générique $\eta$, $X$ un schéma séparé, plat et de type fini sur  $S$, $\ell$ un nombre premier
inversible sur $S$, $Z \subset X$ un fermé rare. Alors il existe une extension finie $\eta'$ de
$\eta$ de degré premier à $\ell$, et une $\ell'$-altération projective $h : \tilde X \rightarrow X$
au-dessus de  $S'\rightarrow S$, où $S'$ est le normalisé de $S$ dans $\eta'$, avec $\tilde X$
régulier et quasi-projectif sur $S'$, un diviseur à croisements normaux stricts $\tilde T$ sur
$\tilde X$,  et une partie fermée finie $\Sigma$ de $S'$ tels que :

(i) en dehors de $\Sigma$, $\tilde X \ra S'$ est lisse et $\tilde T \ra S'$ est un diviseur relatif
à croisements normaux ;

(ii) localement pour la topologie étale autour de chaque point géométrique $x$ de $\tilde X_{s'}$,
où $s' = \SP \,k'$ appartient à $\Sigma$, le couple $(\tilde X, \tilde T)$ est isomorphe au couple
formé de $$ X'=S'[t_1,\cdots,t_n,u_1^{\pm 1}, \cdots, u_s^{\pm
1}]/(t_1^{a_1}...t_r^{a_r}u_1^{b_1}\cdots u_s^{b_s}-\pi), $$ $$ T' = V(t_{r+1} \cdots t_m) \subset
X' $$ autour du point $(u_i = 1)$, $(t_j = 0)$, avec $1 \le r \le m \le n$, pour des entiers $> 0$
$a_1, \cdots, a_r, b_1, \cdots, b_s$ tels que $\mathrm{pgcd}(p,a_1,\cdots,a_r,b_1,\cdots, b_s) = 1$,
$p$ désignant l'exposant caractéristique de $k'$, et $\pi$ une uniformisante locale en $s'$ ;

(iii) $h^{-1}(Z)_{\mathrm{red}}$ est un sous-diviseur de $\cup_{s' \in \Sigma} (\tilde X_{s'}) \cup
\tilde T$.  \end{theoreme}

Le théorème d'uniformisation \ref{0-uniformisation l'} a le corollaire suivant, dit théorème
d'uniformisation « faible », où n'apparaît plus de nombre premier $\ell$ :
\begin{corollaire}\label{0-uniformisation faible} Soient $X$ un schéma noethérien quasi-excellent,
$Z$ un fermé rare de $X$. Il existe une famille finie de morphismes $(p_i : X_i \ra X)_{i \in I}$,
couvrante pour la topologie des altérations et telle que, pour tout $i \in I$ :

(i) $X_i$ soit régulier et connexe,

(ii) $p_i^{-1}(Z)$ soit le support d'un diviseur à croisements normaux stricts.  \end{corollaire} La
topologie des altérations est définie de manière analogue à celle des $\ell'$-altérations. Elle est
plus fine que la topologie étale, et aucune contrainte n'est imposée sur le degré des extensions
résiduelles génériques, \textit{cf.} \refext{II}{2-alt}. On peut démontrer \ref{0-uniformisation
faible} indépendamment de \ref{0-uniformisation l'}, en suivant seulement les étapes (A) et (B)
décrites plus haut, à l'aide, dans (A), d'une forme faible du théorème de Cohen-Gabber. Le théorème
de modification de (C) est inutile, il n'y a pas besoin de faire appel à la résolution canonique des
singularités en caractéristique nulle. La résolution des singularités toriques de Kato-Niziol
suffit. Cette démonstration est exposée dans VII.

\medskip La démonstration de \ref{0-finitude1} est donnée dans (\refext{XIII}{13-section
constructibilité forte}), en application de \ref{0-uniformisation l'} et du théorème de pureté
cohomologique absolu. L'énoncé de \ref{0-finitude1} comporte deux assertions :

(i) la constructibilité des $R^qf_*F$ pour tout $q$,

(ii) l'existence d'un entier $N$ (dépendant de $(f,F)$) tel que $R^qf_*F = 0$ pour $q \ge N$,
autrement dit, le fait que $Rf_*$ envoie $D^b_c(X,\ZZ/n\ZZ)$ dans $D^b_c(Y,\ZZ/n\ZZ)$, l'indice $c$
désignant la sous-catégorie pleine de $D^b$ formée des complexes à cohomologie constructible. 

Dans (\refext{XIII}{13-section constructibilité forte}), ces deux assertions sont démontrées
simultanément. On peut toutefois démontrer (i) en n'invoquant que le théorème d'uniformisation
faible \ref{0-uniformisation faible} (et le théorème de pureté). Ceci est fait dans
(\refext{XIII}{13-constructibilité sans annulation}). L'idée est la suivante. Si, dans
\ref{0-uniformisation faible}, les morphismes $p_i$ étaient propres, on pourrait, après réduction au
cas où $f$ est une immersion ouverte et $F$ un faisceau constant, se ramener, par descente
cohomologique propre, au cas de l'immersion du complément d'un diviseur à croisements normaux
stricts dans un schéma régulier, justiciable du théorème de pureté. Cependant, les $p_i$ ne sont pas
propres en général. Gabber tourne la difficulté à l'aide du théorème de constructibilité générique
de Deligne \sgaquatreetdemi{Th. finitude}{1.9 (i)} et d'un théorème de descente cohomologique
« orientée », établi dans XII. Ce théorème utilise la notion de \textit{produit orienté de topos},
due à Deligne \cite{AG_Tokyo-Kyoto@Laumon}. Les définitions et propriétés de base sont rappelées
dans XI. Un résultat clé est un théorème de changement de
base pour des « tubes » (\refext{XI}{11-ch. base}).

D'après des exemples classiques de Nagata, un schéma noethérien quasi-excellent n'est pas
nécessairement de dimension finie. Si $Y$ est de dimension finie, alors $Rf_*$ est de dimension
cohomologique finie, d'après un théorème de Gabber, exposé dans XVIII, et par suite (ii) découle de
(i). Gabber prouve, plus précisément, que si $X$ est un schéma noethérien, strictement local, de
dimension $d >0$, et $\ell$ un nombre premier inversible sur $X$, alors, pour tout ouvert $U$ de
$X$, on a $\mathrm{cd}_{\ell}(U) \le 2d-1$.

L'hypothèse de quasi-excellence dans \ref{0-finitude1} est essentielle, comme le montre l'exemple
donné dans XIX, d'une immersion ouverte $j : U \ra X$ de schémas noethériens, avec 2 inversible sur
$X$, telle que $R^1j_*(\ZZ/2\ZZ)$ ne soit pas constructible. 

Gabber a démontré des variantes de \ref{0-finitude1} pour les faisceaux d'ensembles ou de groupes
non commutatifs \cite{Conference-Illusie@Gabber}~:

\begin{theoreme}\label{0-finitude2}
Soit $f : X \ra Y$ un morphisme de type fini entre schémas noethériens. Alors :

(1) Pour tout faisceau d'ensembles constructible $F$ sur $X$, $f_*F$ est constructible.

(2) Si $Y$ est quasi-excellent, et si $L$ est un ensemble de nombres premiers inversibles sur $Y$,
pour tout faisceau de groupes constructible et de $L$-torsion $F$ sur $X$, $R^1f_*F$ est
constructible.
\end{theoreme}

La démonstration est donnée dans XXI. Elle ne fait pas appel aux
théorèmes d'uniformisation précédents. Elle utilise, pour le point clé, une technique
d'ultra-produits, et un théorème de rigidité pour les coefficients non abéliens, dû également à
Gabber, qui est établi dans XX. Ce théorème est une variante d'un théorème de Fujiwara-Gabber pour
les coefficients abéliens \cite[6.6.4, 7.1.1]{TUBULAR@Fujiwara} (signalée dans (\textit{loc. cit.},
6.6.5). Il s'énonce ainsi :

\begin{theoreme}\label{0-rigidité}
Soient $(X,Y)$ un couple hensélien,
où $X = \SP A$, $Y = V(I)$, l'idéal $I$ étant supposé de type fini, $\widehat X$ le complété
$I$-adique de $X$, $U$ un ouvert de $X$ contenant $X-Y$, et $\widehat U = \widehat X \times_X U$.
Alors, pour tout champ en groupoïdes ind-fini $C$ sur $U$, la flèche de restriction $\Gamma(U,C) \ra
\Gamma(\widehat U, C \vert \widehat U)$ est une équivalence.
\end{theoreme}
La démonstration donnée dans XXI est indépendante de \cite{TUBULAR@Fujiwara}.

\medskip Le reste de la partie II est consacré à trois autres applications des théorèmes
d'uniformisation de la partie I.

(a) \textit{Lefschetz affine}. Il s'agit de généralisations des théorèmes de \sga{4}{xiv}{3.1} (pour
les morphismes affines entre schémas de type fini sur un corps) et \sga{4}{xix}{6.1} (pour les
morphismes affines de type fini de schémas excellents de caractéristique nulle), ainsi que du
théorème de Gabber pour les morphismes affines de schémas de type fini sur un trait
\cite{PV@Illusie}. L'énoncé principal est le suivant (\refext{XV}{15-theoreme-principal}) :

\begin{theoreme}\label{0-Lefschetz affine}
Soit $f : X \ra Y$ un morphisme affine de type fini, où
$Y$ est un schéma noethérien quasi-excellent, muni d'une fonction de dimension $\delta_Y$, et soit
$n$ un entier inversible sur $Y$. Alors, pour tout faisceau constructible $F$ de $\ZZ/n\ZZ$-modules
sur $X$, et tout entier $q$, on a :
\begin{equation}\label{0-Lefschetz affine1}
\delta_Y(R^qf_*F) \le \delta_X(F) - q
\end{equation}
\end{theoreme}
Une \textit{fonction de dimension} $\delta$ sur un
schéma $T$ est une fonction $\delta : T  \ra \ZZ$ telle que $\delta(y) = \delta(x) - 1$ si $y$ est
une spécialisation étale immédiate de $x$. Cette notion est due à Gabber.  Elle généralise celle de
\textit{dimension rectifiée} introduite dans \sgacourt{4}{xiv}{}. Elle est définie et étudiée dans XIV.
Dans \ref{0-Lefschetz affine1}, la fonction $\delta_X$ est reliée à $\delta_Y$ par $\delta_X(x) =
\delta_Y(f(x)) + \mathrm{degtr}(k(x)/k(f(x))$, et  pour un faisceau $G$ sur $X$ (resp.  $Y$)
$\delta_X(G)$ (resp. $\delta_Y(G)$) désigne la borne supérieure des $\delta_X(x)$ (resp.
$\delta_Y(x)$) pour $G_x \ne 0$. La démonstration de \ref{0-Lefschetz affine} se fait par réduction
au théorème de Gabber cité plus haut, pour les schémas de type fini sur un trait, à l'aide du
théorème d'uniformisation « faible » \ref{0-uniformisation faible} et du théorème de descente
cohomologique orientée de XII. De \ref{0-Lefschetz affine} résulte :
\begin{corollaire}\label{0-Lefschetz affine2} Si $X$ est local noethérien, hensélien,
quasi-excellent, et de dimension $d$, de corps résiduel $k$, et si $U$ est un ouvert affine de $X$,
alors, pour tout nombre premier $\ell$ inversible sur $X$, on a \begin{equation}\label{0-Lefschetz
affine3} \mathrm{cd}_{\ell}(U) \le d + \mathrm{cd}_{\ell}(k).  \end{equation} \end{corollaire} En
particulier, si $X$ est intègre, de corps des fractions $K$, on en déduit $\mathrm{cd}_{\ell}(K) = d
+ \mathrm{cd}_{\ell}(k)$, formule conjecturée dans \sga{4}{x}{3.1}. Les valeurs possibles de
$\mathrm{cd}_{\ell}(U)$, pour $U$ ouvert, non nécessairement affine, de $X$ sont étudiées dans
XVIII. Gabber donne également des contre-exemples à \ref{0-Lefschetz affine3} lorsqu'on omet
l'hypothèse de quasi-excellence. 

\medskip (b) \textit{Une nouvelle démonstration de la conjecture de pureté absolue}. La
démonstration de cette conjecture qui est donnée dans \cite{AZUMINO@Fujiwara} utilise, dans sa
dernière partie, des techniques de $K$-théorie algébrique (résultats de Thomason). Gabber a annoncé
dans \cite{Conference-Deligne@Gabber} qu'on peut éviter tout recours à la $K$-théorie algébrique, en
utilisant, à la place, la forme raffinée du théorème de de Jong \ref{0-deJong l'}{(2)}. Cette
nouvelle démonstration est exposée en détail dans XVI. Ce chapitre contient en outre une théorie de
classes fondamentales généralisées (due à Gabber), utilisée pour construire une théorie de
morphismes de Gysin pour les morphismes d'intersection complète lissifiables, généralisant les
constructions de \sgaquatreetdemicourt{Cycle}{}.

\medskip (c) \textit{Complexes dualisants}. La notion de complexe dualisant est due à Grothendieck.
L'unicité, l'existence et les propriétés générales des complexes dualisants sont étudiées dans
\sgacourt{5}{I}{}. Toutefois, dans (\textit{loc.  cit.}) l'existence n'est établie qu'en caractéristique
nulle, ou sous des hypothèses d'existence de résolution des singularités, et de validité du théorème
de pureté absolue (conjecturale à l'époque). Dans le cas de schémas de type fini sur un schéma
régulier de dimension $\le 1$, l'existence est prouvée inconditionnellement par Deligne dans
\sgaquatreetdemicourt{Dualité}{}. Dans le cas général, l'existence, et la théorie de dualité locale qui
en résulte, ont été annoncées par Gabber dans \cite{Conference-Deligne@Gabber}. Le chapitre XVII
expose cette théorie en détail. Si $X$ est un schéma noethérien, et $\Lambda = \ZZ/n\ZZ$, où $n$ est
un entier inversible sur $X$, un \textit{complexe dualisant} sur $X$ est un objet de
$D^b_c(X,\Lambda)$ tel que le foncteur $D_K = R\mathscr{H}om(-,K)$ envoie $D^b_c(X,\Lambda)$ dans
$D^b_c(X,\Lambda)$ et que, pour tout $L \in D^b_c(X,\Lambda)$, la flèche de bidualité $L \ra
D_KD_K(L)$ soit un isomorphisme. Cette définition diffère légèrement de celle de \sga{5}{I}{1.7},
voir (\refext{XVII}{17-definition-complexe-dualisant}).  L'unicité, à décalage, et torsion près par
un $\Lambda$-module inversible, est prouvée dans \sga{5}{I}{2.1}. Le résultat principal de XVII est
que, si $X$ est excellent et admet une fonction de dimension, au sens précisé plus haut, alors $X$
admet un complexe dualisant. De plus, ces complexes dualisants ont les propriétés de fonctorialité
attendues, et, si $X$ est régulier, le faisceau constant $\Lambda_X$ est dualisant. Cette dernière
assertion était une conjecture dans \sgacourt{5}{I}{}, démontrée dans le cas de caractéristique nulle.
Nous renvoyons le lecteur à l'introduction de XVII pour des énoncés plus complets, et des
indications sur la méthode de démonstration, dont les ingrédients essentiels sont le théorème de
finitude \ref{0-finitude1} et le théorème d'algébrisation partielle de V (voir (A) \textit{supra}).

\renewcommand\thesubsection{\textbf{\thesection.\arabic{subsection}}}

\ifx\danslelivre\undefined
\bibliography{../fichiers-configuration-latex/bibliographie}
\bibliographystyle{../fichiers-configuration-latex/smfalpha}
\end{document}
\fi
